\section{Conclusion and Future Work}

\towrite[Aim: open the chapter in one short paragraph. Cover: that it brings the work together. Introduce no new results or sources.]

\subsection{Introduction}

\towrite[Aim: outline the chapter. Cover: restate the research question and the order of the sections that follow.]

\subsection{Summary of the Research}

\towrite[Aim: retell the work briefly. Cover: the problem; what was built; how it was tested; the main findings.]

\subsection{Main Contributions}

\towrite[Aim: state what the work delivered. Cover: the actual outcomes against the expected contributions in Section 1.10, including the framework, protocol, evidence and any data released.]

\subsection{Extent to Which Objectives Were Achieved}

\towrite[Aim: close the loop on the objectives. Cover: each objective from Section 1.7, marked as met, partly met or not met, with a short justification; a direct answer to the research question.]

\subsection{Limitations of the Study}

\towrite[Aim: state the limits on the conclusions. Cover: a brief summary of Section 6.9 and how far the findings can be generalised.]

\subsection{Recommendations for Future Research}

\towrite[Aim: suggest what other researchers should do next. Cover: larger and more diverse samples; custom classifiers; other paradigms or devices; less controlled environments; long-term studies. Base each on your findings.]

\subsection{Final Remarks}

\towrite[Aim: end the dissertation in one short paragraph. Cover: the wider significance of the work and the main message for the reader.]
